\section*{Erd\H{o}s Problem 702: forbidding singleton intersections}

\paragraph{Statement qualification and claim.}
For each fixed \(k\geq4\), and for all sufficiently large \(n\)
depending on \(k\), a \(k\)-uniform family on \([n]\) with more
than \(\binom{n-2}{k-2}\) members contains two sets intersecting
in exactly one point. The dependence on sufficiently large
\(n\) is essential. It appears explicitly in Frankl's primary
paper, although absent from the displayed statement at
\url{https://www.erdosproblems.com/702}.
We state the established theorem as an external input, check
the sharp example, and give explicit counterexamples to the
unqualified variants.

\paragraph{The known theorem, with its hypotheses retained.}
Frankl proved that for each integer \(k\geq4\) there is
\(n_0(k)\) such that whenever \(n>n_0(k)\) and
\(\mathcal F\subseteq\binom{[n]}k\) satisfies
\[
 |A\cap B|\ne1\qquad(A,B\in\mathcal F,\ A\ne B),
\]
one has
\[
 |\mathcal F|\leq\binom{n-2}{k-2}.                              \tag{1}
\]
This is the principal external theorem used here:
P. Frankl, \emph{On families of finite sets no two of which
intersect in a singleton}, Bull. Austral. Math. Soc. 17 (1977),
125--134,
\url{https://doi.org/10.1017/S0004972700025521}.
Taking the contrapositive of (1) proves the correctly quantified
assertion. This use does not purport to supply a new proof of
Frankl's structural theorem.

\paragraph{Sharpness of the bound.}
Fix two points \(a,b\in[n]\) and take
\[
 \mathcal F=\{A\in\binom{[n]}k:\{a,b\}\subseteq A\}.
\]
There are exactly \(\binom{n-2}{k-2}\) choices. Every two
members meet in at least \(a,b\), so none intersect in a
singleton. Thus the strict inequality in the question is
necessary, and the bound (1) is sharp.

\paragraph{Why the omitted size qualification cannot be dropped.}
For any \(k\geq4\), take \(n=k+1\) and all \(k\)-subsets of
\([k+1]\). The family has \(k+1\) members, whereas
\[
 \binom{n-2}{k-2}=\binom{k-1}{k-2}=k-1.
\]
Two different members intersect in \(k-1\geq3\) elements.
This disproves the literal assertion for arbitrary \(n\);
it does not contradict Frankl's theorem.

\paragraph{Why the restriction \(k\geq4\) also matters.}
For \(k=3\) and \(n=4t\), partition the ground set into \(t\)
blocks of size four and take every triple inside each block.
There are \(4t=n\) triples. Two triples from the same block
intersect in two elements, and triples from different blocks
are disjoint. Thus no singleton intersection occurs, although
\[
 n>\binom{n-2}{1}=n-2.
\]
This gives arbitrarily large counterexamples when \(k=3\).
For \(k=2\), a matching of two or more edges already exceeds
\(\binom{n-2}{0}=1\) and has no singleton intersection.

\paragraph{A direct link-family reduction.}
For each \(x\in[n]\), define
\(\mathcal F_x=\{A\setminus\{x\}:x\in A\in\mathcal F\}\).
The original family has no singleton intersection if and
only if every \(\mathcal F_x\) is intersecting.
Indeed, two disjoint members of \(\mathcal F_x\) lift to sets
whose intersection is exactly \(\{x\}\), and conversely.
For \(n-1\geq2(k-1)\), the usual Erd\H{o}s--Ko--Rado theorem
on each link gives
\[
 |\mathcal F_x|\leq\binom{n-2}{k-2},\qquad
 k|\mathcal F|\leq n\binom{n-2}{k-2}.
\]
This last bound explains why treating the links separately
does not recover the sharp result: Frankl's theorem controls
their compatibility and removes the extra factor \(n/k\).

\paragraph{Scope.}
The literal statement without a condition on \(n\) is disproved
by the explicit example above. The intended large-\(n\)
statement follows completely from the precisely stated
published theorem; its threshold and its two necessary
qualifications have been checked directly.

\paragraph{Completion estimate.}
\textbf{COMPLETION: 100\%}
